\pagestyle{fancy}\thispagestyle{plain}

\rule{6.8in}{1pt}\vskip.1in

{\large Assignment 1: \csname topic01\endcsname}\smallskip

\rule{6.8in}{1pt}\vskip.1in

\textit{Due by \HWdueTime,   eastern, on \dueDay, \csname dateWeek01\endcsname}

\rule{6.8in}{1pt}\vskip.1in

{\bf\textit{Suggested readings  for this problem set:}} Chapter 1, except for proof by contradiction.
% \begin{itemize}[itemsep=1pt]
% \item Chapter 1, except for proof by contradiction.
% \end{itemize}
\smallskip

All readings are from \bookAuthors, \emph{\bookName}.

\rule{6.8in}{1pt}\vskip.1in

{\bf\textit{Extra Practice Problems}} from \emph{\bookName} (not to turn in):

\begin{itemize}[itemsep=1pt]
\item With answers or hints (in the back of the book):
 \begin{itemize}
 \item Section 1.1, \#1(adgj), 2(adji), 3(adgi), 5(ad), 6(a)
 \item Section 1.2, \#2(ac), 4(ac), 5(ad), 7(a), 10(a), 11(a), 12(a)
 \item Section 1.3, \#1(ad), 3(a), 5(ac), 7(ac)
 \item Section 1.4, \#1, 4(a), 6(a), 8, 12(ab), 15(a)

\end{itemize}

  \item \href{https://dmzb.github.io/teaching/2024Spring220/handouts/220-H01-logic.pdf}{Handout 1}
\end{itemize}

\rule{6.8in}{1pt}\vskip.1in

 {\bf\textit{Assignment:}} due \dueDay, \csname dateWeek01\endcsname, \HWdueTime, via Gradescope (\gradescopeCode):

\rule{6.8in}{1pt}\vskip.1in

\begin{enumerate}
% \item Add boolean algebra problem, truth table problem (for contrapositive); ``which make mathematical sense'', see V for more

\item Write the negation of each of the following statements.

 \begin{enumerate}
\item All triangles are isosceles.
\item Every door in the building was locked.
\item Some even numbers are multiples of three.
\item Every real number is less than 100.
\item Every integer is positive or negative.
\item If $f$ is a polynomial function, then $f$ is continuous at $0$.
\item If $x^2 > 0$, then $x > 0$.
\item There exists a $y \in \mathbf{R}$ such that $xy = 1$.
\item $(2 > 1)$ and $(\forall x, x^2 > 0)$
\item $\forall \epsilon > 0$, $\exists \delta > 0$ such that if $|x| < \delta$, then $|f(x)| < \epsilon$.
\end{enumerate}

\item Write the converse, contrapositive, and negation of each of the following implications.

 \begin{enumerate}
\item If a quadrilateral is a rectangle, then it has two pairs of parallel sides.
\item $(P \wedge \neg Q) \Rightarrow R$
\item $P  \Rightarrow (R \Rightarrow \forall x, Q(x))$
\end{enumerate}

\item Let $P$ and $Q$ be statements. Write the truth table for
  \begin{enumerate}
  \item $(\neg P) \vee Q$
  \item $(P \wedge (\neg Q)) \Rightarrow Q$
  \end{enumerate}

  
\item Are the statements $(P \vee Q) \wedge R$ and $P \vee (Q \wedge R)$ equivalent? If so, give a proof. If not, explain why by giving a counterexample.

  (Two statement forms are equivalent if they have the same truth tables, and here, a counterexample simply means some choice of truth values for $P,Q,$ and $R$ such that the two statement forms give different outputs.)
\item Let $P$ and $Q$ be statements.
  \begin{enumerate}
  \item Prove that $\neg(P \Rightarrow Q)$ is equivalent to $P \wedge \neg Q$.
  \item Prove that $\neg(P \Rightarrow Q)$ is \emph{not} equivalent to $\neg P \wedge Q$.
  \item Give an example of statements $P$ and $Q$ such that $\neg P \Rightarrow \neg Q$ is true and $\neg(P \Rightarrow Q)$ is false.
  \end{enumerate}

% \item \textbf{TODO} Write a problem where they prove that two statement forms are equivalent using properties rather than truth tables. 

% \item Let $n$ and $m$ be integers. Suppose that $nm$ is even. Is it true that both $n$ and $m$ are even? If so, give a proof. If not, give a counterexample (i.e., give an example that demonstrates that this is false.)

\item Suppose that $n$ is an even integer, and let $m$ be any integer. Prove that $nm$ is even.   
  
\item Suppose that $n$ is an odd integer. Prove that $n^2$ is an odd integer. (Hint: an integer $n$ is odd if and only if there exists an integer $k$ such that $n = 2k+1$.)

\item Prove that if $n^2$ is even, then $n$ is even.
  (Click \href{https://en.wikipedia.org/wiki/Proof_by_contrapositive}{{\color{red}here}} for a hint. Also: you are allowed to use the results of previous problems.)  
  \end{enumerate}



%%% Local Variables: 
%%% mode: latex
%%% TeX-master: "../assignments-S24-220"
%%% End: 